\documentclass[8pt,letterpaper]{extarticle}
\usepackage[margin=0.45in]{geometry}
\usepackage{multicol}
\usepackage{amsmath}
\usepackage{amssymb}
\usepackage{enumitem}
\usepackage{titlesec}
\usepackage{booktabs}
\usepackage{array}
\usepackage{microtype}
\usepackage{parskip}
\usepackage{xcolor}
\usepackage{mdframed}
\usepackage{fancyhdr}

% Compact list settings
\setlist[itemize]{noitemsep, topsep=1pt, parsep=0pt, leftmargin=10pt}
\setlist[enumerate]{noitemsep, topsep=1pt, parsep=0pt, leftmargin=12pt}

% Section formatting
\titleformat{\section}{\bfseries\small\uppercase}{}{0pt}{}[\titlerule]
\titlespacing{\section}{0pt}{4pt}{2pt}
\titleformat{\subsection}{\bfseries\footnotesize}{}{0pt}{}
\titlespacing{\subsection}{0pt}{3pt}{1pt}

% Code/formula box
\newmdenv[
  backgroundcolor=gray!10,
  linecolor=gray!40,
  linewidth=0.5pt,
  innerleftmargin=4pt,
  innerrightmargin=4pt,
  innertopmargin=2pt,
  innerbottommargin=2pt,
  skipabove=2pt,
  skipbelow=2pt
]{formulabox}

\pagestyle{fancy}
\fancyhf{}
\renewcommand{\headrulewidth}{0.4pt}
\fancyhead[L]{\textbf{CSDS 340 -- Exam 3 Cheat Sheet} \quad {\small Lec 14--21}}
\fancyhead[R]{\small Exam: 4/20 $\cdot$ 60 min $\cdot$ Calculator OK}
\fancyfoot[C]{\thepage}

\setlength{\columnsep}{8pt}
\setlength{\parindent}{0pt}

\begin{document}
\begin{multicols}{3}

%=============================================
\section{Ensemble Methods (Lec 14)}
%=============================================

\subsection{Core Idea}
Combine multiple classifiers to reduce generalization error. Can use same/different models, data, or hyperparameters.

\textbf{Majority Voting:}
\begin{formulabox}
$\hat{y} = \text{mode}(C_1(x), C_2(x), \ldots, C_m(x))$
\end{formulabox}

\textbf{Weighted Voting (prob outputs):}
\begin{formulabox}
$\text{Score}_\text{class} = \sum_k w_k \cdot p_{k,\text{class}}$\\
$\hat{y} = \arg\max(\text{Score}_\text{class})$
\end{formulabox}

\subsection{Bagging}
\begin{itemize}
  \item Train each classifier on a \textbf{randomly sampled subset with replacement}
  \item All classifiers trained independently (\textbf{parallel})
  \item Reduces variance; different views of same data
\end{itemize}

\subsection{AdaBoost}
Weak learner = \textbf{decision stump} (1-node tree). Labels: $-1$ and $+1$. Sequentially learns from mistakes.

\begin{formulabox}
\texttt{Init: } $w_i = 1/n$ \quad (uniform, $\sum w_i = 1$)\\
\texttt{For } $j = 1\ldots m$:\\
\quad $C_j = \text{train}(X, y, w)$\\
\quad $\hat{y} = \text{predict}(C_j, X)$\\
\quad $\varepsilon = w \cdot (\hat{y} \neq y)$ \quad [weighted error]\\
\quad $\alpha_j = 0.5 \cdot \ln\!\left(\frac{1-\varepsilon}{\varepsilon}\right)$\\
\quad $w_i \leftarrow w_i \cdot e^{-\alpha_j \hat{y}_i y_i}\ \forall i$\\
\quad Normalize $w$ to sum to 1\\
\texttt{Final: } $\hat{y} = \text{sign}\!\left(\sum_j \alpha_j \cdot \text{predict}(C_j, X)\right)$
\end{formulabox}

\textbf{Weight update logic:}
\begin{itemize}
  \item Correct: $\hat{y} \cdot y = +1 \Rightarrow e^{-\alpha} \Rightarrow$ weight \textbf{decreases}
  \item Wrong: $\hat{y} \cdot y = -1 \Rightarrow e^{+\alpha} \Rightarrow$ weight \textbf{increases}
  \item Large $\alpha_j$ $\Rightarrow$ classifier trusted more in final vote
\end{itemize}

%=============================================
\section{Gradient Boosting (Lec 15)}
%=============================================

\begin{center}
\begin{tabular}{@{}lll@{}}
\toprule
 & \textbf{AdaBoost} & \textbf{Grad. Boost} \\
\midrule
Weak learner & Stump & DT (3--4 lvls) \\
Learns from  & Weighted err & Gradient of loss \\
Loss         & Exp. loss & Logistic loss \\
\bottomrule
\end{tabular}
\end{center}

\begin{formulabox}
$L_i = -[y_i\log p_i + (1-y_i)\log(1-p_i)]$\\
$p = \text{sigmoid of predicted log-odds}$\\[2pt]
\textbf{Step 0:} $\hat{y}_0 = \log(\#\text{pos}/\#\text{neg})$,\; $p_0 = \#\text{pos}/n$\\[2pt]
\textbf{Each round } $m$:\\
\quad a) Residuals: $r_{im} = y_i - p_i$\\
\quad b) Fit new tree to $(X, \text{residuals})$\\
\quad c) Each leaf $R_{jm}$: $\gamma_{jm} = \frac{\sum(y_i-p_i)}{\sum p_i(1-p_i)}$\\
\quad d) Update: $F_m(x) = F_{m-1}(x) + \eta\cdot\gamma_{jm}$\\[2pt]
\textbf{Inference:} $\hat{y}_\text{new} = \hat{y}_0 + \eta(\gamma_{j_1}+\gamma_{j_2}+\cdots+\gamma_{j_m})$
\end{formulabox}

\textbf{Worked Example} (2 pos, 1 neg):
\begin{formulabox}
$\hat{y}_0 = \ln(2/1) = 0.69,\quad p_0 = 0.67$\\
$r(y=1) = 1-0.67 = 0.33$;\quad $r(y=0) = -0.67$\\
$\gamma_{\text{leaf}(y=1)} = \frac{0.33+0.33}{0.67{\times}0.33{\times}2} = 1.49$\\
$\gamma_{\text{leaf}(y=0)} = \frac{-0.67}{0.67{\times}0.33} = -3.12$
\end{formulabox}
Trees are \textbf{sequential}; stop when residuals $<$ threshold or $M$ trees reached.

%=============================================
\section{Linear Regression -- OLS (Lec 16)}
%=============================================

\subsection{Definition}
Below is the definition of linear regression, the learning rules and model decision parameters

\begin{formulabox}
$\hat{y} = w^\top x + b$ \quad ($y$ is continuous)\\[2pt]
\textbf{Loss (MSE):} $L(w,b) = \frac{1}{n}\sum_i (y_i - \hat{y}^{(i)})^2$\\[2pt]
$\frac{\partial L}{\partial w_j} = -\frac{2}{n}\sum_i(y_i-\hat{y}^{(i)})x_j^{(i)}$\\
$\frac{\partial L}{\partial b} = -\frac{2}{n}\sum_i(y_i-\hat{y}^{(i)})$
\end{formulabox}

\textbf{Regularization}
\begin{formulabox}
Ridge (L2): $\text{MSE} + \lambda\|w\|_2^2$ \quad (small, distributed $w$)\\
LASSO (L1): $\text{MSE} + \lambda\|w\|_1$ \quad (sparse, some $w=0$)\\
Elastic Net: $\text{MSE} + \lambda_1\|w\|_1 + \lambda_2\|w\|_2^2$
\end{formulabox}

\textbf{Residual Plots:} plot $\varepsilon = y - \hat{y}$. Good: symmetric around 0. Bad: skewed/patterned.

\textbf{Pearson Correlation (Feature Selection)}
\begin{formulabox}
$r = \frac{\sigma_{w_j y}}{\sigma_{w_j}\cdot\sigma_y}$\\
$\sigma_{w_j y} = \sum_i(x_j^{(i)}-\mu_j)(y^{(i)}-\mu_y)$\\
$\sigma_{w_j} = \sqrt{\sum_i(x_j^{(i)}-\mu_j)^2}$\\
Range: $-1$ to $+1$;\; $|r|\approx 1$ = strong linear relationship
\end{formulabox}
\textbf{Multi-collinearity:} features are linearly dependent ($r \approx \pm 1$ between features). Fix: drop one.

\textbf{Evaluation Metrics}
\begin{formulabox}
$\text{MSE} = \frac{1}{n}\sum(y_i-\hat{y}_i)^2$\\
$\text{MAE} = \frac{1}{n}\sum|y_i-\hat{y}_i|$\\
$R^2 = 1 - \frac{\text{SSE}}{\text{SST}} = 1 - \frac{\text{MSE}}{\text{Var}(y)}$\\
Train: $0 \leq R^2 \leq 1$;\quad Test: $R^2$ can be \textbf{negative}
\end{formulabox}

%=============================================
\section{Polynomial \& Random Forest Reg (Lec 17)}
%=============================================

\subsection{Polynomial Regression}
Add polynomial features, then train OLS. Degree $d$: add all products up to degree $d$.
\begin{formulabox}
4 features, degree 2 $\Rightarrow$ 15 total features:\\
$1, x_1,x_2,x_3,x_4, x_1^2,x_2^2,x_3^2,x_4^2,$\\
$x_1x_2, x_1x_3, x_1x_4, x_2x_3, x_2x_4, x_3x_4$
\end{formulabox}

\subsection{Decision Tree Regression}
Same greedy split but use \textbf{MSE}; node prediction = \textbf{mean} of samples.
\begin{formulabox}
$I(D) = \frac{1}{N_D}\sum_{i\in D}(y_i - \bar{y}_D)^2$\\
$\text{Gain}(x_i) = I(D) - \frac{N_l}{N}I(D_l) - \frac{N_r}{N}I(D_r)$
\end{formulabox}

\subsection{Random Forest Regression}
Ensemble of DT Regressors using \textbf{bagging}; final prediction = \textbf{average} of all trees.

\begin{center}\footnotesize
\begin{tabular}{@{}llll@{}}
\toprule
\textbf{Model} & \textbf{Linear?} & \textbf{Loss}\\
\midrule
OLS & Yes & MSE\\
Ridge & Yes & MSE$+\lambda\|w\|_2^2$\\
LASSO & Yes & MSE$+\lambda\|w\|_1$\\
Polynomial & No & MSE on poly feats\\
Rand. Forest & No & MSE per node\\
\bottomrule
\end{tabular}
\end{center}

\begin{center}\footnotesize
\begin{tabular}{@{}lll@{}}
\toprule
 & \textbf{Bagging} & \textbf{Boosting}\\
\midrule
Training & Parallel & Sequential\\
Sampling & Random w/ replace. & Weighted\\
Learners & Independent & Each fixes prev.\\
Reduces & Variance & Bias\\
\bottomrule
\end{tabular}
\end{center}

%=============================================
\section{K-Means Clustering (Lec 17--18)}
%=============================================

Unsupervised; no labels $y$. Find $K$ natural clusters minimizing within-cluster SSE.
\begin{formulabox}
$D = \sum_i\sum_k r_{ik}\|x^{(i)}-\mu_k\|^2$\\
$r_{ik} = 1$ if sample $i$ in cluster $k$, else $0$\\
$\mu_k = \frac{\sum_i r_{ik}\, x^{(i)}}{\sum_i r_{ik}}$ \quad [cluster mean]
\end{formulabox}

\textbf{Algorithm:}
\begin{enumerate}
  \item Pick $K$ random points as initial $\mu_k$
  \item Assign each sample to nearest centroid
  \item Update $\mu_k$ = mean of assigned samples
  \item Repeat until convergence
\end{enumerate}

\textbf{Elbow Method (choosing $K$):} Run K-means for $k=1\ldots K'$; plot $D$ vs $k$; choose $k$ at the ``elbow.'' $D=0$ is trivial (K=N).

\textbf{Silhouette Value}
\begin{formulabox}
$a^{(i)} = $ avg dist to points in \emph{same} cluster\\
$b^{(i)} = $ avg dist to points in \emph{nearest other} cluster\\
$s^{(i)} = \frac{b-a}{\max(b,a)}$\\[2pt]
$s \approx +1$: well-matched; \quad $s \approx 0$: boundary; \quad $s \approx -1$: wrong cluster
\end{formulabox}

%=============================================
\section{Hierarchical Clustering (Lec 19)}
%=============================================

No need to choose $K$ upfront; builds a \textbf{dendrogram}.

\textbf{Agglomerative (bottom-up):}
\begin{enumerate}
  \item Start: $N$ clusters (one per sample)
  \item Compute pairwise distances
  \item Merge closest pair
  \item Repeat until 1 cluster
\end{enumerate}
\begin{formulabox}
Single linkage: $\min\, d(u,v)$\\
Complete linkage: $\max\, d(u,v)$
\end{formulabox}
Cut dendrogram at desired level to extract $K$ clusters. Application: phylogenetic trees.

%=============================================
\section{Neural Networks (Lec 19--21)}
%=============================================

\subsection{Architecture}
Input Layer $\to$ Hidden Layers $\to$ Output Layer. Each neuron: $a = \sigma(w^\top x + b)$.

\textbf{Counting Parameters:}
\begin{formulabox}
Weights $= \sum(\text{dim}_\text{prev} \times \text{dim}_\text{next})$\\
\textit{E.g.} input=2, hidden=[3,3], output=1:\\
Weights $= 2{\times}3 + 3{\times}3 + 3{\times}1 = 18$;\quad Neurons $= 9$\\
Weight label: $w_{ij}$ = from neuron $i$ to neuron $j$
\end{formulabox}

\textbf{Activation Functions}
\begin{formulabox}
Sigmoid: $\sigma(z) = \frac{1}{1+e^{-z}}$;\quad $\frac{d\sigma}{dz} = \sigma(z)(1-\sigma(z))$\\
ReLU: $\max(0,z)$ \quad [preferred in deep nets]\\
Step function: $\mathbf{1}[z>0]$ \quad [Perceptron]
\end{formulabox}

\subsection{Forward Propagation}
Layer by layer, left $\to$ right: $a_n = \sigma\!\left(\sum w_{\text{prev},n}\cdot a_\text{prev} + b_n\right)$

\textbf{Example} (2$\to$[3]$\to$1, all $w=1$, $b=0$):
\begin{formulabox}
$x=(0,0)$: $a_3{=}\sigma(0){=}0.5$,\; $a_5{=}\sigma(1.0){=}0.73$\\
$x=(0,1)$: $a_3{=}0.73$,\; $a_5{=}\sigma(1.46){=}0.81$\\
$x=(1,1)$: $a_3{=}0.73$,\; $a_5{=}\sigma(1.46){=}0.85$
\end{formulabox}

\textbf{Loss (MSE):} $L = \frac{1}{n}\sum(y_i - \hat{y}_i)^2$;\quad $\frac{dL}{da_5} = -2(y - a_5)$

\subsection{Backpropagation (Chain Rule)}
Compute \textbf{right-to-left}; reuse intermediate gradients.
\begin{formulabox}
$\frac{dL}{dw_{35}} = \frac{dL}{da_5}\cdot\frac{da_5}{dw_{35}}$\\
$\frac{da_5}{dw_{35}} = a_5(1-a_5)\cdot a_3$\\[3pt]
$\frac{dL}{dw_{13}} = \frac{dL}{da_5}\cdot\frac{da_5}{da_3}\cdot\frac{da_3}{dw_{13}}$\\
$\frac{da_5}{da_3} = a_5(1-a_5)\cdot w_{35}$;\quad $\frac{da_3}{dw_{13}} = a_3(1-a_3)\cdot x_1$\\[3pt]
Multi-path (neuron 3 $\to$ 6,7,8): sum gradients over all paths
\end{formulabox}

\textbf{Training Loop:}
\begin{enumerate}
  \item Init weights (small random)
  \item Forward prop $\to$ compute $\hat{y}$
  \item Compute loss $L(\hat{y}, y)$
  \item Backward prop $\to$ gradients
  \item Update: $w \leftarrow w - \eta\cdot\frac{dL}{dw}$
  \item Repeat until convergence
\end{enumerate}

%=============================================
\section{Deep Learning (Lec 21)}
%=============================================

ANN with $>3$ layers = Deep Learning. Sigmoid in hidden layers $\to$ \textbf{vanishing gradients}; fix with \textbf{ReLU} in hidden layers.
\begin{itemize}
  \item More layers = more abstract feature representations
  \item Deeper nets require careful weight initialization to train stably
  \item Use \textbf{cross-entropy loss} for multi-class classification
  \item Optimization via \textbf{SGD} (stochastic gradient descent) or variants (Adam, etc.)
  \item Training mode: gradients enabled; Eval mode: gradients disabled
\end{itemize}

%=============================================
\section{Quick Facts \& Exam Traps}
%=============================================

\subsection{Ensemble}
\begin{itemize}
  \item AdaBoost labels are $-1$ and $+1$, \textbf{NOT} 0 and 1
  \item Sample weights $\neq$ model weights (as in Adaline)
  \item $\varepsilon$ = weighted error; $\alpha_j$ larger when $\varepsilon$ small (better classifier)
  \item Correct pred: weight $\downarrow$;\quad Wrong pred: weight $\uparrow$
  \item Bagging = parallel;\quad Boosting = sequential
  \item Gradient boosting residual $r_i = y_i - p_i$ (not $\hat{y}$)
  \item $\gamma_{jm} = \sum r / \sum p(1-p)$ for each leaf
  \item Inference: sum all $\gamma$ contributions $\times \eta$
\end{itemize}

\subsection{Regression}
\begin{itemize}
  \item OLS = Adaline but $y$ continuous (no threshold)
  \item Pearson $r$ measures \textbf{linear} correlation only
  \item $R^2$ on test can be \textbf{negative} (worse than mean predictor)
  \item $R^2{=}1$ = perfect fit; $R^2{=}0$ = always predicts $\mu_y$
  \item LASSO $\to$ sparse weights (feature selection)
  \item Ridge $\to$ small distributed weights
  \item DT Regression: node prediction = \textbf{mean}; split by MSE gain
  \item Random Forest Regression: bagging + MSE splits
\end{itemize}

\subsection{Clustering \& Neural Networks}
\begin{itemize}
  \item K-means: $D=0$ requires $K=N$ (trivial/useless)
  \item Silhouette: $+1{=}$good, $0{=}$boundary, $-1{=}$wrong cluster
  \item Hierarchical: no need to pick $K$ first; cut dendrogram
  \item Single linkage = min dist;\quad Complete = max dist
  \item NN weights per layer pair $= d_\text{prev} \times d_\text{next}$
  \item Backprop: \textbf{right-to-left}; reuse intermediate gradients
  \item Multiple paths $\to$ \textbf{sum} gradients from all paths
  \item $\frac{d\sigma}{dz} = \sigma(z)(1-\sigma(z))$ \quad $\leftarrow$ memorize!
  \item ReLU fixes vanishing gradient in deep nets
\end{itemize}

\end{multicols}
\end{document}
